% 第 4 章 §4.2.3（印刷页 179–184；PDF p.189–194）：Hitchin–Kobayashi 对应；自命题 4.2.16 证明的收尾起，至定理 4.2.17 证明的第 1 步止
% 对应 工程文件/work/en/ch04.txt 第 796–1200 行（印刷页 179–184）
% 本块编号公式：(4.20)–(4.28)，共 9 条；(4.29)、(4.30) 与定理 4.2.17 证明的第 2–4 步归 ch04d
% 本块定理类环境：定理 4.2.17、注 4.2.18、定义 4.2.19、引理 4.2.20、命题 4.2.21、引理 4.2.22、注 4.2.23、
%   引理 4.2.24、引理 4.2.25、引理 4.2.26、命题 4.2.27（编号由 main.tex 自动生成，不手写）
% 记号：$\mathcal{E}$ 为全纯丛（底本作花体 E），$E$ 为其底光滑丛；$\mathcal{S},\mathcal{Q},\mathcal{M},\mathcal{N}$ 亦为花体；
%   $\bar\partial_A$ 为连接 $A$ 的 $(0,1)$-部分；$\|\cdot\|$ 为 $L^2$ 范数，$|\cdot|$ 为逐点范数；
%   「Chern 连接」沿用 ch04a 的译法（术语表作「Chern 连接」）。
% 块界说明：承接 ch04b 末（命题 4.2.16 的证明在 p.178 页末中断，按其接缝说明另起「（续）」proof 环境），
%   止于定理 4.2.17 证明第 1 步末（第 2–4 步见 ch04d）。

\begin{proof}[命题 4.2.16 的证明（续）]
以及
\begin{equation}
\operatorname{tr}\left(\pi\sqrt{-1}\Lambda[\Phi,\Phi^*]\right)=|\pi\Phi(I-\pi)|^2.
\label{eq:4.20}
\tag{4.20}
\end{equation}
把 (4.17) 代入 (4.4)，并利用 (4.20)，得
\[
\deg\mathcal{E}=\operatorname{rank}(\mathcal{S})\mu(\mathcal{E})-\frac{1}{2\pi}\left(\|\pi\Phi(I-\pi)\|^2+\|\beta\|^2\right).
\]
这就证明了 $\mu(\mathcal{S})\le\mu(\mathcal{E})$。并且，等号成立当且仅当上式右端的两项都为零；也就是说，全纯结构与 Higgs 场分裂。
\end{proof}

本节要证明的主要结果是命题 4.2.16 的逆命题。

\begin{theorem}[Hitchin {[35]}，Simpson {[57]}]
若 $(\mathcal{E},\Phi)$ 是多稳定的，则它容许一个满足 (4.17) 的度量。
\end{theorem}

\begin{remark}
线丛 $\mathcal{L}$ 的情形是直接的。事实上，秩为 $1$ 时 $[\Phi,\Phi^*]$ 这一项为零，故 (4.17) 归结为在 $L$ 上找一个常曲率度量（constant curvature metric）。设 $H$ 为任一度量，令 $H_\varphi=e^\varphi H$，其中 $\varphi$ 是函数。则
\[
F_{(\bar\partial_L,H_\varphi)}=F_{(\bar\partial_L,H)}+\bar\partial\partial\varphi,
\]
于是只要找到 $\varphi$ 使
\[
\Delta\varphi=2\sqrt{-1}\Lambda(F_{(\bar\partial_L,H)})-2\deg(L).
\]
问题便解决了。由 Hodge 定理，该方程有解的唯一条件是右端的积分为零（参看 [26, p. 84]），而由 (4.3) 它确实为零。
\end{remark}

为了在高秩情形证明定理 4.2.17，构造近似临界度量将是重要的。设
$0\subset(\mathcal{E}_1,\Phi_1)\subset\cdots\subset(\mathcal{E}_\ell,\Phi_\ell)=(\mathcal{E},\Phi)$
是 Higgs 丛 $(\mathcal{E},\Phi)$ 的 Harder–Narasimhan 滤过。令 $\mathcal{Q}_i=\mathcal{E}_i/\mathcal{E}_{i-1}$，$\mu_i=\mu(\mathcal{Q}_i)$。则存在光滑分裂 $E=\bigoplus_i\mathcal{Q}_i$，并且给定埃尔米特度量 $H$ 后，可使该分裂正交。于是存在一个良定义的自同态
\begin{equation}
\mu_{(\mathrm{Gr}(\mathcal{E},\Phi),H)}=
\begin{pmatrix}
\mu_1& & \\
& \ddots& \\
& &\mu_\ell
\end{pmatrix},
\label{eq:4.21}
\tag{4.21}
\end{equation}
其中块 $\mu_i$ 的维数为 $\operatorname{rank}\mathcal{Q}_i$。

\begin{definition}
若
\[
\sup\left|f_{((\bar\partial_{\mathcal{E}},H),\Phi)}-\mu_{(\mathrm{Gr}(\mathcal{E},\Phi),H)}\right|<\varepsilon,
\]
则称 $(\mathcal{E},\Phi)$ 上的一个度量为 $\varepsilon$-近似临界（$\varepsilon$-approximate critical）的。
\end{definition}

注意，$\mathcal{E}$ 的 $\bar\partial$-算子关于该分裂可以写成上三角形式，而严格上三角部分由扩张类决定。用一个块对角的复规范变换作用，可使扩张类变得任意小。此外，若丛 $\mathcal{Q}_i$ 连同其诱导的 Higgs 场容许 Hermitian–Yang–Mills–Higgs 联络，则可以把这些度量相加，得到下面的结果（更多细节见 [14]）。

\begin{lemma}
设 $(\mathcal{E},\Phi)$ 是秩为 $n$ 的不稳定 Higgs 丛，并设定理 4.2.17 对秩小于 $n$ 的 Higgs 丛已经证明。则对任意 $\varepsilon>0$，$(\mathcal{E},\Phi)$ 上都存在 $\varepsilon$-近似临界度量。
\end{lemma}

\paragraph{预备估计}
回忆映射 (4.15)。一个关键点是下面的先验估计（a priori estimate）。

\begin{proposition}
设 $(\mathcal{E},\Phi)$ 是 Higgs 丛。存在只依赖于 $X$ 与 $E$ 上的度量以及 $\|h[\mathcal{E},\Phi]\|$ 的常数 $C_1,C_2>0$，使得
\[
\sup|\Phi|^2\le C_1+C_2\sup\left|\sqrt{-1}\Lambda(F_A+[\Phi,\Phi^*])\right|.
\]
\end{proposition}

需要下面的引理。

\begin{lemma}[参看 {[59, p. 27]}]
对矩阵 $P$，存在只依赖于 $P$ 的特征值的常数 $C_1,C_2>0$，使得
\[
|[P,P^*]|^2\ge C_1|P|^4-C_2(1+|P|^2).
\]
\end{lemma}

\begin{proof}
取一组酉基使 $P=S+N$，其中 $S$ 对角而 $N$ 严格上三角。由假设 $|S|$ 有界。容易看出，此时只需证明存在 $C>0$，使得对一切严格上三角矩阵 $N$ 都有 $|[N,N^*]|\ge C|N|^2$。若不然，则通过缩放可以找到序列 $N_j$，$|N_j|=1$，而 $[N_j,N_j^*]\to0$。取子列后可设 $N_j\to N$，且 $[N,N^*]=0$、$|N|=1$。但这是矛盾的。事实上，若 $a_1,\dots,a_n$ 与 $b_1,\dots,b_n$ 分别是 $N$ 的行与列，则从 $NN^*=N^*N$ 的对角元可以读出 $|a_i|^2=|b_i|^2$，$i=1,\dots,n$。但 $b_1=0$，由这一等式推出 $a_1=0$。这又给出 $b_2=0$，从而 $a_2=0$。如此继续，最终得到 $N=0$，矛盾。
\end{proof}

还需要下面的计算：
\begin{equation}
\begin{aligned}
[[P,P^*],P]&=(PP^*-P^*P)P-P(PP^*-P^*P)\\
&=2PP^*P-P^*P^2-P^2P^*,\\
\langle[[P,P^*],P],P\rangle&=\operatorname{tr}\bigl([[P,P^*],P]P^*\bigr)=\operatorname{tr}\bigl((2PP^*P-P^*P^2-P^2P^*)P^*\bigr)\\
&=2\operatorname{tr}(PP^*)^2-2\operatorname{tr}(P^2(P^*)^2),\\
\langle\operatorname{ad}([P,P^*])P,P\rangle&=|[P,P^*]|^2.
\end{aligned}
\label{eq:4.22}
\tag{4.22}
\end{equation}

\begin{proof}[命题 4.2.21 的证明]
把 $\Phi$ 看作 $\mathcal{K}\otimes\operatorname{End}\mathcal{E}$ 的全纯截面。我们还要用到三个简单事实。第一，若 $H$ 是 $E$ 上的埃尔米特度量、$\tilde{H}$ 是它在 $\operatorname{End}E$ 上诱导的度量，则 $F_{(\operatorname{End}\mathcal{E},\tilde{H})}=\operatorname{ad}F_{(\mathcal{E},H)}$，其中 $\operatorname{ad}$ 表示曲率自同态通过交换子作用。第二，若 $\tilde{H}$、$h$ 分别是 $\operatorname{End}\mathcal{E}$ 与 $\mathcal{K}$ 上的埃尔米特度量，则
\begin{equation}
F_{(\mathcal{K}\otimes\operatorname{End}\mathcal{E},h\otimes\tilde{H})}=F_{(\operatorname{End}\mathcal{E},\tilde{H})}+F_{(\mathcal{K},h)}\cdot I.
\label{eq:4.23}
\tag{4.23}
\end{equation}
第三，若 $s$ 是带酉联络 $A$ 与曲率 $F_A$ 的向量丛的一个全纯截面，则
\[
\begin{aligned}
\Delta|s|^2&=-2\bar\partial^*\bar\partial|s|^2=2\sqrt{-1}\Lambda\partial\bar\partial|s|^2=2\sqrt{-1}\Lambda\partial\langle s,d'_As\rangle\\
&=2\sqrt{-1}\Lambda\langle d'_As,d'_As\rangle+2\sqrt{-1}\Lambda\langle s,d''_Ad'_As\rangle\\
&=2|d'_As|^2+2\sqrt{-1}\Lambda\langle s,F_As\rangle\\
&=2|d_As|^2-2\langle\sqrt{-1}\Lambda F_As,s\rangle,
\end{aligned}
\]
这给出 Weitzenböck 公式（Weitzenböck formula）
\begin{equation}
\Delta|s|^2=2|d_As|^2-2\langle\sqrt{-1}\Lambda F_As,s\rangle.
\label{eq:4.24}
\tag{4.24}
\end{equation}
现在利用 (4.22)、(4.23)、(4.24) 以及引理 4.2.22，得
\[
\begin{aligned}
\Delta|\Phi|^2&\ge-2\langle\sqrt{-1}\Lambda F_{(\mathcal{K}\otimes\operatorname{End}\mathcal{E},h\otimes\tilde{H})}\Phi,\Phi\rangle\\
&\ge-2\langle\sqrt{-1}\Lambda F_{(\operatorname{End}\mathcal{E},\tilde{H})}\Phi,\Phi\rangle-C_3|\Phi|^2\\
&=-2\langle\operatorname{ad}(\sqrt{-1}\Lambda F_{(\mathcal{E},H)})\Phi,\Phi\rangle-C_3|\Phi|^2\\
&=2\langle\operatorname{ad}(\sqrt{-1}\Lambda[\Phi,\Phi^*])\Phi,\Phi\rangle-2\langle\operatorname{ad}(\sqrt{-1}\Lambda(F_{(\mathcal{E},H)}+[\Phi,\Phi^*]))\Phi,\Phi\rangle-C_3|\Phi|^2\\
&\ge C_1|\Phi|^4-C_2(1+|\Phi|^2)-C_4\sup\left|\sqrt{-1}\Lambda(F_{(\mathcal{E},H)}+[\Phi,\Phi^*])\right||\Phi|^2.
\end{aligned}
\]
最后，在 $|\Phi|^2$ 的最大值处左端非正。由于 $C_1>0$，命题立即得证。
\end{proof}

\begin{remark}
注意 (4.16) 中的符号对这一论证至关重要（参看 [37]）。
\end{remark}

最后，存在性的证明将基于 Donaldson 在 [18] 中的漂亮论证。这需要引入泛函 $J=J(A,\Phi)$，其定义如下。对 $E$ 的埃尔米特自同态 $\varphi$，令
\[
\nu(\varphi)=\sum_{i=1}^n|\lambda_i|,\qquad N^2(\varphi)=\int_X\nu^2(\varphi)\,\frac{\omega}{2\pi},
\]
其中 $\lambda_i$ 是 $\varphi$ 的（逐点）特征值。然后定义
\begin{equation}
J(A,\Phi)=N(f_{(A,\Phi)}-\mu(\mathcal{E})).
\label{eq:4.25}
\tag{4.25}
\end{equation}

下面我们证明 Donaldson 的两个结果（见 [18] 的引理 2 与引理 3），这里把它们改写成 Higgs 丛的情形。

\begin{lemma}
设 $(A,\Phi)$ 是底丛为 $\mathcal{E}$ 的 Higgs 丛。设它构成 Higgs 丛的扩张
\[
0\longrightarrow\mathcal{M}\longrightarrow\mathcal{E}\longrightarrow\mathcal{N}\longrightarrow0,
\]
并且 $\mu(\mathcal{N})\le\mu(\mathcal{E})\le\mu(\mathcal{M})$。则
\[
(\operatorname{rank}\mathcal{M})(\mu(\mathcal{M})-\mu(\mathcal{E}))+(\operatorname{rank}\mathcal{N})(\mu(\mathcal{E})-\mu(\mathcal{N}))\le J(A,\Phi).
\]
\end{lemma}

\begin{proof}
关于正交分裂 $E=M\oplus N$，并记 $F_E$、$F_M$、$F_N$ 为 $(\mathcal{E},H)$ 的Chern 连接的曲率及其诱导曲率，有
\[
\sqrt{-1}\Lambda F_E=
\begin{pmatrix}
\sqrt{-1}\Lambda F_M+b_M&-(d''_A)^*\beta\\
-((d''_A)^*\beta)^*&\sqrt{-1}\Lambda F_N+b_N
\end{pmatrix},
\]
其中 $\beta$ 是第二基本形式，且 $b_M=-\sqrt{-1}\Lambda(\beta\wedge\beta^*)$、$b_N=-\sqrt{-1}\Lambda(\beta^*\wedge\beta)$。注意 $\operatorname{tr}b_M=-\operatorname{tr}b_N=|\beta|^2$。类似地，若写
\[
\Phi=\begin{pmatrix}\Phi_M&\varphi\\0&\Phi_N\end{pmatrix},
\]
则
\[
[\Phi,\Phi^*]=
\begin{pmatrix}
[\Phi_M,\Phi^*_M]+\varphi\wedge\varphi^*&\varphi\wedge\Phi^*_N+\Phi^*_M\wedge\varphi\\
\Phi_N\wedge\varphi^*+\varphi^*\wedge\Phi_M&[\Phi_N,\Phi^*_N]+\varphi^*\wedge\varphi
\end{pmatrix}.
\]
由此得到
\[
f_{(A,\Phi)}=
\begin{pmatrix}
f_M+b_M+\sqrt{-1}\Lambda\varphi\wedge\varphi^*&\cdots\\
\cdots&f_N+b_N+\sqrt{-1}\Lambda\varphi^*\wedge\varphi
\end{pmatrix}.
\]
于是（参看 [18, p. 271]）
\[
\begin{aligned}
\nu(f_{(A,\Phi)}-\mu(\mathcal{E}))\ge{}&\left|\operatorname{tr}(\sqrt{-1}\Lambda F_M)-(\operatorname{rank}\mathcal{M})\mu(\mathcal{E})+|\beta|^2+|\varphi|^2\right|\\
&+\left|\operatorname{tr}(\sqrt{-1}\Lambda F_N)-(\operatorname{rank}\mathcal{N})\mu(\mathcal{E})-|\beta|^2-|\varphi|^2\right|
\end{aligned}
\]
从而
\[
\begin{aligned}
J(A,\Phi)&\ge\int_X\nu(f_{(A,\Phi)}-\mu(\mathcal{E}))\,\frac{\omega}{2\pi}\\
&\ge\left|\int_X\left(\operatorname{tr}(\sqrt{-1}\Lambda F_M)-(\operatorname{rank}\mathcal{M})\mu(\mathcal{E})+|\beta|^2+|\varphi|^2\right)\frac{\omega}{2\pi}\right|\\
&\quad+\left|\int_X\left(\operatorname{tr}(\sqrt{-1}\Lambda F_N)-(\operatorname{rank}\mathcal{N})\mu(\mathcal{E})-|\beta|^2-|\varphi|^2\right)\frac{\omega}{2\pi}\right|\\
&\ge(\operatorname{rank}\mathcal{M})(\mu(\mathcal{M})-\mu(\mathcal{E}))+(\operatorname{rank}\mathcal{N})(\mu(\mathcal{E})-\mu(\mathcal{N})).
\end{aligned}
\]
\end{proof}

\begin{lemma}
设 $(A_0,\Phi_0)$ 是秩为 $n$ 的稳定 Higgs 丛，它构成 Higgs 丛的扩张
$0\to\mathcal{S}\to\mathcal{E}\to\mathcal{Q}\to0$。假设定理 4.2.17 对秩小于 $n$ 的 Higgs 丛已经证明。则可以在 $(A_0,\Phi_0)$ 的复规范轨道中选取一点 $(A,\Phi)$，使得
\[
J(A,\Phi)<(\operatorname{rank}\mathcal{S})(\mu(\mathcal{E})-\mu(\mathcal{S}))+(\operatorname{rank}\mathcal{Q})(\mu(\mathcal{Q})-\mu(\mathcal{E})).
\]
\end{lemma}

\begin{proof}
首先考虑 $(\mathcal{S},\Phi_{\mathcal{S}})$ 与 $(\mathcal{Q},\Phi_{\mathcal{Q}})$ 的 Harder–Narasimhan 滤过。由引理 4.2.20，可设对任意 $\varepsilon>0$，$\mathcal{S}$ 上存在一个度量使
\[
\sup\left|f_{((\bar\partial_{\mathcal{S}},H_{\mathcal{S}}),\Phi_{\mathcal{S}})}-\mu_{(\mathrm{Gr}(\mathcal{S},\Phi_{\mathcal{S}}),H_{\mathcal{S}})}\right|<\varepsilon,
\]
对 $\mathcal{Q}$ 亦然。给 $E=S\oplus Q$ 配上这两个度量之和。这对应于轨道中 $(A_0,\Phi_0)$ 的某个点 $(A,\Phi)$。其次，由于 $(A_0,\Phi_0)$（从而是 $(A,\Phi)$）是单的，可以进一步设
\[
-\bar\partial^*_{A_0}\beta+\sqrt{-1}\Lambda\left(\varphi\wedge\Phi^*_{\mathcal{Q}}+\Phi^*_{\mathcal{S}}\wedge\varphi\right)=0;
\]
见 (4.8)。这可以由形如 $g=\begin{pmatrix}1&\varphi\\0&1\end{pmatrix}$ 的复规范变换实现。特别地，$S$ 与 $Q$ 上的 $\bar\partial$-算子保持不变，故近似临界结构仍然成立。在此基础上，再作一次规范变换，使 $(A,\Phi)$ 与 $(A_0,\Phi_0)$ 相同，只是 $\beta$ 与 $\varphi$ 乘上了 $t$。此时 $f_{(A,\Phi)}-\mu(\mathcal{E})$ 是块对角的，其元素为
\begin{equation}
\begin{aligned}
&f_{\mathcal{S}}-\mu_{(\mathrm{Gr}(\mathcal{S},\Phi_{\mathcal{S}}),H_{\mathcal{S}})}+\mu_{(\mathrm{Gr}(\mathcal{S},\Phi_{\mathcal{S}}),H_{\mathcal{S}})}-\mu(\mathcal{E})+t^2\left(b_{\mathcal{S}}+\sqrt{-1}\Lambda\varphi\wedge\varphi^*\right),\\
&f_{\mathcal{Q}}-\mu_{(\mathrm{Gr}(\mathcal{Q},\Phi_{\mathcal{Q}}),H_{\mathcal{Q}})}+\mu_{(\mathrm{Gr}(\mathcal{Q},\Phi_{\mathcal{Q}}),H_{\mathcal{Q}})}-\mu(\mathcal{E})+t^2\left(b_{\mathcal{Q}}+\sqrt{-1}\Lambda\varphi^*\wedge\varphi\right).
\end{aligned}
\label{eq:4.26}
\tag{4.26}
\end{equation}
由于 $(\mathcal{E},\Phi)$ 稳定，$\mu(\mathcal{E})$ 严格大于 $\mathcal{S}$ 的任一子层的最大斜率，并严格小于 $\mathcal{Q}$ 的任一商层的最小斜率。这说明当 $t$ 与 $\varepsilon$ 取得足够小时，(4.26) 的第一行负定而第二行正定。由此得到
\[
\begin{aligned}
\nu(f_{(A,\Phi)}-\mu(\mathcal{E}))\le{}&(\operatorname{rank}\mathcal{S})(\mu(\mathcal{E})-\mu(\mathcal{S}))+(\operatorname{rank}\mathcal{Q})(\mu(\mathcal{Q})-\mu(\mathcal{E}))\\
&-2t^2\left(|\beta|^2+|\varphi|^2\right)+\mathrm{O}(\varepsilon).
\end{aligned}
\]
不妨设 $\|\beta\|^2+\|\varphi\|^2=1$。按 [18] 的论证，还可以设 $|\beta|$、$|\varphi|$ 关于 $\varepsilon$ 一致有界。现在固定 $t$ 并取 $\varepsilon$ 足够小，即得结论。
\end{proof}

\paragraph{存在性定理}
下一节引入 Yang–Mills–Higgs 流之后，我们将证明下面的结果。

\begin{lemma}
在任一复规范轨道中都存在序列 $(A_i,\Phi_i)$，满足以下条件：
\begin{enumerate}[label=(\roman*), leftmargin=2.6em, itemsep=0.2em]
\item $(A_i,\Phi_i)$ 使 $J$ 极小化；
\item 若 $f_{(A_j,\Phi_j)}=\sqrt{-1}\Lambda(F_{A_j}+[\Phi_j,\Phi_j^*])$，则 $\sup|f_{(A_j,\Phi_j)}|$ 关于 $j$ 一致有界；
\item $\|d_{A_j}f_{(A_j,\Phi_j)}\|_{L^2}\to0$，且 $\|[f_{(A_j,\Phi_j)},\Phi_i]\|_{L^2}\to0$。
\end{enumerate}
\end{lemma}

\begin{annote}
底本此引理中下标 $i$、$j$ 混用：(i) 用 $(A_i,\Phi_i)$，(ii)(iii) 改用 $j$，而 (iii) 末尾又出现 $\Phi_i$。按上下文，整列序列应统一记作 $(A_j,\Phi_j)$，(iii) 中的 $\Phi_i$ 显系 $\Phi_j$ 之误。译文照排底本，不改动下标。[来源：上下文推断]（参见疑点清单 Q089）
\end{annote}

接下来要用到规范理论中最基本的结果之一，这里按黎曼面的情形叙述。

\begin{proposition}[Uhlenbeck，{[64]}]
固定 $p\ge2$。设 $\{A_j\}$ 是一列 $L^p_1$-联络，且 $\|F_{A_j}\|_{L^p}$ 一致有界。则存在一列酉规范变换 $g_j\in L^p_2$ 与一个光滑酉联络 $A_\infty$，使得在取适当的子列之后，$g_j(A_j)\to A_\infty$ 在 $L^p_1$ 中弱收敛、在 $L^p$ 中强收敛。
\end{proposition}

有了这些结果，现在来证明存在性定理。

\begin{proof}[定理 4.2.17 的证明]
显然只需假设 $(\mathcal{E},\Phi)$ 稳定。此外，由注 4.2.18，可以按归纳法进行。设结论对一切秩小于 $n=\operatorname{rank}E$ 的丛都已证明。

\emph{第 1 步}（极限丛 $(\mathcal{E}_\infty,\Phi_\infty)$）。按引理 4.2.26 取 $J$ 的极小化序列（minimizing sequence）。由于该序列位于同一个复规范轨道中，Hitchin 映射的像 $h[A_i,\Phi_i]$ 不变。于是由命题 4.2.21，$\Phi_i$ 一致有界。再由引理 4.2.26 (ii)，这又推出对任意 $p$，$\|F_{A_j}\|_{L^p}$ 有界。因此由命题 4.2.27，可以设存在光滑联络 $A_\infty$，使得若写 $\bar\partial_{A_j}=\bar\partial_{A_\infty}+a_j$，则 $a_j\to0$ 在 $L^p_1$ 中弱收敛。由 Sobolev 嵌入定理，特别地可设 $a_j\to0$ 于某个 $C^\alpha$。于是 $F_{A_j}\to F_{A_\infty}$ 在 $L^p$ 中弱收敛。由全纯性条件，
\[
0=\bar\partial_{A_j}\Phi_j=\bar\partial_{A_\infty}\Phi_j+[a_j,\Phi_j].
\]
$\bar\partial_{A_\infty}$ 的椭圆正则性给出估计 $\|\Phi_j\|_{L^2_1}\le C\|\Phi_j\|_{L^2}$。对前一式求导得
\begin{equation}
\bar\partial^*_{A_\infty}\bar\partial_{A_\infty}\Phi_j+\bar\partial^*_{A_\infty}[a_j,\Phi_j]=0.
\label{eq:4.27}
\tag{4.27}
\end{equation}
由 Cauchy–Schwarz 不等式与上述估计，有
\begin{equation}
\|\bar\partial^*_{A_\infty}[a_j,\Phi_j]\|_{L^2}\le C_1\|a_j\|_{L^4_1}\|\Phi_j\|_{L^4}+C_2\|\Phi_j\|_{L^2}.
\label{eq:4.28}
\tag{4.28}
\end{equation}
现在可以设 $\{a_j\}$ 在 $L^4_1$ 中有界，并利用 Laplacian $\bar\partial^*_{A_\infty}\bar\partial_{A_\infty}$ 的椭圆正则性以及包含关系 $L^2_1\hookrightarrow L^4$、$L^2_2\hookrightarrow C^\alpha$，由 (4.27) 与 (4.28) 得到估计 $\|\Phi_j\|_{C^\alpha}\le C\|\Phi_j\|_{L^2}$。由于 $\Phi_j$ 一致有界，其 $L^2$ 范数有界，故可设 $\Phi_j$ 在 $C^\alpha$ 中收敛到某个 $\Phi_\infty$。此外，由 $\Phi_j$ 的全纯性，可以写
\[
\bar\partial_{A_\infty}\Phi_\infty=\bar\partial_{A_\infty}(\Phi_\infty-\Phi_j)-[a_j,\Phi_j],
\]
而 $[a_j,\Phi_j]\to0$ 于 $C^\alpha$，故 $\bar\partial_{A_\infty}\Phi_\infty=0$ 弱成立。于是由 Weyl 引理，$\Phi_\infty$ 实际上全纯，从而 $(\mathcal{E}_\infty,\Phi_\infty)$ 是 Higgs 丛。
\end{proof}
